\section{Global Systemic Stagnation Hypothesis (GSSH)}
\subsection{Definition}
GSSH posits that the world economy has entered a structurally stagnant state caused by the interaction of four constraint layers: (1) diminishing marginal impact of monetary policy, (2) demographic aging and fertility collapse, (3) AI-driven productivity concentration in capital-intensive sectors, and (4) geopolitical fragmentation that reduces trade and capital mobility.

\subsection{Mechanism}
The feedback lock forms when weak monetary transmission fails to lift demand, demographic drag limits labor responsiveness, technological gains accrue narrowly, and fragmentation raises transaction frictions. The resulting loop suppresses productivity, investment, and political stability.

\subsection{Indicators and Signals}
Observed signals include weakening fiscal and monetary multipliers, dependency ratio acceleration, divergence of risk premiums across blocs, asset inflation decoupled from real activity, and early financial triggers such as stablecoin peg stress, UST auction yield spikes, and CRE refinancing failures.

\subsection{Comparative Cases}
Japan's post-1990 stagnation (demographic dominance), the euro area's post-sovereign constraints (monetary/geopolitical interplay), and recent trade bifurcation episodes illustrate partial manifestations of the GSSH loop.

\subsection{Reference Diagram}
\placeholderfigure{Layer\_Interaction.pdf}{GSSH constraint layers and systemic lock}
